\subsection{Real-world validation study}
\label{sec:realworld}

The curated corpus (Section~\ref{sec:evaluation}) was designed to contain
representative defects, so its defect rates measure detection capability rather
than prevalence. To estimate \emph{base rates} on independently authored
workflows, and to measure extractor fidelity on code the authors did not write,
we mined public GitHub repositories across all four supported frameworks.

\paragraph{Mining pipeline.}
Using GitHub code search we identified candidate repositories for each
framework, then streamed each repository through clone, extraction, and cleanup
(so disk never accumulates all clones at once). This yielded 930 workflow graphs
from 253 distinct repositories: LangGraph~402, AutoGen~361, CrewAI~114, and
Google~ADK~53. Extraction used the static AST fallback
(\texttt{scripts/ast\_extractor.py}); the runtime extractors of
Section~\ref{sec:graph_model} read the compiled graph object and are expected to
be more faithful, but require per-repository dependency installation that is
infeasible at this scale. Mined workflows are small: a median of 2~non-sentinel
nodes (mean~3.9), predominantly static linear pipelines.

\paragraph{Validation protocol.}
On a stratified 120-workflow sample (40~LangGraph, 20~CrewAI, 35~AutoGen,
25~ADK) we ran two independent LLM-agent passes. First, an agent reconstructed a
\emph{ground-truth} graph from source alone (blind to the extractor output),
following the extractor's node-identifier conventions so the two align; we then
computed node/edge precision--recall and node-kind accuracy against the AST
extraction. Second, for each of the 187 flagged defects, one agent labelled it
\textsc{real}, \textsc{extraction-artifact}, or \textsc{intentional} from
source, and an adversarial verifier independently attempted to overturn each
label, with disagreements recorded as \textsc{arguable}.

\begin{table}[t]
\centering
\small
\caption{AST-extractor fidelity on real GitHub workflows ($n=120$ ground-truth
graphs). Stub-based validation (Section~\ref{sec:accuracy}) reported perfect
precision/recall; on real code fidelity is substantially lower and strongly
framework-dependent---ADK, whose topology is most implicit, is the worst case.}
\label{tab:realworld_fidelity}
\begin{tabular}{@{}lrrrr@{}}
\toprule
\textbf{Framework} & \textbf{$n$} & \textbf{Edge P} & \textbf{Edge R} & \textbf{Kind acc} \\
\midrule
LangGraph & 40  & 0.956 & 0.682 & 0.647 \\
CrewAI    & 20  & 0.700 & 0.692 & 0.929 \\
AutoGen   & 35  & 0.957 & 0.770 & 0.900 \\
ADK       & 25  & 0.403 & 0.367 & 0.943 \\
\midrule
\textbf{Overall} & \textbf{120} & \textbf{0.798} & \textbf{0.644} & \textbf{0.829} \\
\bottomrule
\end{tabular}
\end{table}

\paragraph{Extractor fidelity.}
Table~\ref{tab:realworld_fidelity} reports the result. Node detection stays
strong (overall precision~0.91, recall~0.85), but edge recall falls to 0.644
(95\% CI $[0.57,0.71]$, bootstrap over workflows) and
node-kind accuracy to 0.829---far below the perfect scores on programmatic
stubs---and edge recall ranges from 0.37 (ADK) to 0.77 (AutoGen). The frameworks
that encode topology \emph{implicitly} (ADK's nested \texttt{SequentialAgent}/
\texttt{ParallelAgent} composition, AutoGen's participant lists with dynamic
speaker selection) are hardest: the extractor must guess an edge set the source
never states explicitly, and for ADK it is wrong more often than right. On
LangGraph, whose edges are explicit, the dominant error is instead node-kind
(\texttt{tool}$\rightarrow$\texttt{llm}: nodes that invoke tools inside their
body without declared bindings).

\paragraph{The fallback, not the ceiling.}
The 0.64 figure is a property of the static AST fallback used for large-scale
mining, not of the runtime extractors (Section~\ref{sec:graph_model}), which read
the compiled framework object and are therefore faithful by construction. On the
repository's own runnable example workflows---the only code we execute, as running
untrusted mined code would require a sandbox we do not deploy---the AST fallback
recovers 0.71 of the runtime extractor's LangGraph edges, closely matching its
0.68 on mined LangGraph code; the edges it misses are chiefly the conditional
targets of \texttt{add\_conditional\_edges} maps that a static reader cannot
resolve. For ADK the two extractors name nodes differently, so a direct edge
recall is not meaningful, but the runtime graph is markedly richer (16 versus
${\sim}5$ edges), consistent with ADK's implicit topology being the AST fallback's
worst case. The higher-fidelity runtime path thus exists but does not scale to
untrusted code; this gap between the faithful-but-unscalable and
scalable-but-lossy extraction paths is precisely why extraction fidelity, rather
than the check algorithms, is the binding constraint.

\begin{table}[t]
\centering
\small
\caption{Adversarial triage of 187 defects flagged on real workflows. Every
structural flag is an extraction artifact; the human-gate policy flag is
predominantly an intentional design choice.}
\label{tab:realworld_triage}
\begin{tabular}{@{}lrrrr@{}}
\toprule
\textbf{Check} & \textbf{Real} & \textbf{Artifact} & \textbf{Intentional} & \textbf{Arguable} \\
\midrule
exit\_reachability     & 0 & 14 & 0  & 0 \\
reverse\_reachability  & 0 & 25 & 0  & 0 \\
dead\_ends             & 0 & 24 & 0  & 1 \\
router\_shape          & 0 & 3  & 0  & 0 \\
tool\_declarations     & 0 & 5  & 0  & 0 \\
human\_presence        & 2 & 9  & 93 & 11 \\
\midrule
\textbf{Total}         & \textbf{2} & \textbf{80} & \textbf{93} & \textbf{12} \\
\bottomrule
\end{tabular}
\end{table}

\paragraph{Defect prevalence.}
Table~\ref{tab:realworld_triage} gives the result: of 187 flagged defects, only
2 (1.1\%, 95\% CI $[0.3,3.8]\%$) are genuine bugs---both missing human
gates---while 43\% are extraction artifacts and 50\% are intentional design
choices. \emph{Every one of the 72 structural flags is an extraction artifact}
(0 genuine; 95\% CI $[0,5.1]\%$), and structural flags are entirely a
LangGraph phenomenon: in the raw study they fire on 78\% of LangGraph workflows
but on 0\% of AutoGen, CrewAI, and ADK workflows, because those extractors build
a connected topology by construction and cannot emit a dead end or unreachable
exit. The structural-defect signal is thus an artifact of the one extractor that
recovers real edges, not a property of the workflows. The \texttt{human\_presence}
policy check fires on 91\% of workflows, but on the validated sample only 2 of
115 firings are genuine (the rest intentional low-risk tasks or artifacts),
reinforcing the context-dependence argued in Section~\ref{sec:evaluation} and
motivating a risk-aware variant.

\paragraph{Risk-aware gating.}
The human-gate check above is deliberately blunt---it asks only whether a
\textsc{human} node exists. A risk-aware variant instead flags a workflow only
when a \emph{sensitive} tool (destructive, write, exec, external-communication,
or financial, by a keyword lexicon grounded in the observed tool names) is
reachable from entry without passing a \textsc{human} node, reusing the sound
coverage check of Section~\ref{sec:verification_methods}. On the curated corpus
this sharpens the signal: flags drop from 10 to 8, declining to flag two
workflows that lack a human node but perform no sensitive action, and pinpointing
the offending operation (e.g.\ \texttt{account\_create}, \texttt{publish}, or
\texttt{auto\_remediate} reachable ungated) rather than merely reporting an absent
gate. On the 930 mined workflows the effect is starker: the naive check fires on
846 (91\%), but the risk-aware check fires on \emph{zero}---none of the mined
workflows declares a sensitive tool, as they are read-only retrieval and
summarisation pipelines that legitimately need no human review. The variant thus
removes the false-positive deluge, but its recall is itself bounded by extraction
fidelity: a sensitive operation invoked inside a node body with no declared tool
binding---the dominant \texttt{tool}$\rightarrow$\texttt{llm} error of
Table~\ref{tab:realworld_fidelity}---is invisible to it. A sound sensitivity
signal therefore depends on the very higher-fidelity extraction this study
identifies as the binding constraint.

\paragraph{Takeaway.}
Topology-level defects are \emph{rare} in publicly available agent workflows,
which are mostly small static pipelines, and the curated-benchmark rates do not
transfer to production prevalence. The practical determinants of static-analysis
utility are therefore extraction fidelity---which varies sharply by framework and
collapses on implicit-topology frameworks like ADK---and policy applicability,
rather than the check algorithms themselves. We read this as redirecting the
research agenda toward higher-fidelity extraction (sandboxed runtime extraction;
AST/runtime cross-validation) and risk-aware policy specification, and toward
giving the static machinery a defect-independent job---pruning provably-inert
runtime monitors---rather than cataloguing topology checks whose real-world
firing is dominated by extraction noise.
